\documentclass[a4paper]{article}
\usepackage{azzam}

\begin{document}

\title{Aljabar Dasar}
\maketitle
\section{Latihan Soal Pemfaktoran dan Manipulasi Aljabar}
\begin{enumerate}
    \item  Nilai dari $\sqrt{5050^2-4950^2}$ adalah \dots

    \item (OSP 2008) Jika $0 < b < a$ dan $a^2+b^2=6ab$, maka nilai $\dfrac{a+b}{a-b}=\dots$
    
    \item Jika $x > 0$ dan $x + \dfrac{1}{x} =  5$, maka nilai $x^3+\dfrac{1}{x^3}$ adalah \dots
    
    \item (OSK 2017) Diketahui $x-y=10$ dan $xy=10$. Nilai $x^4+y^4$ adalah \dots
    
    \item (OSK 2018) Diketahui $x$ dan $y$ bilangan prima dengan $x < y$, dan $x^3+y^3+2018=30y^2-300y+3018$. Nilai $x$ yang memenuhi adalah \dots
    
    \item Jika $a+b+c=0$ untuk suatu bilangan riil $a,b,c$, buktikan bahwa $a^3+b^3+c^3=3abc$.
    
    \item Jika $x=2021^3-2019^3$, maka nilai $\sqrt{\dfrac{x-2}{6}}$ adalah \dots

    \item (AIME 1987)
    Tentukan nilai sederhana dari $\dfrac{(10^4+324)(22^4+324)(34^4+324)(46^4+324)(58^4+324)}{(4^4+324)(16^4+324)(28^4+324)(40^4+324)(52^4+324)}$

    \item (OSK 2017) Jika $\dfrac{(a-b)(c-d)}{(b-c)(d-a)}=-\dfrac{4}{7}$, maka nilai dari $\dfrac{(a-c)(b-d)}{(a-b)(c-d)}$ adalah \dots

    \item (OSK 2019) Diketahui $a+2b=1$, $b+2c=2$, dan $b \neq 0$. Jika $a+nb+2018c = 2019$ maka nilai $n$ adalah \dots

    \item (OSK 2019) Misalkan $a = 2\sqrt{2} - \sqrt{8-4\sqrt{2}}$ dan $b = 2\sqrt{2} + \sqrt{8-4\sqrt{2}}$. Jika $\dfrac{a}{b}+\dfrac{b}{a}=x+y\sqrt{2}$ dengan $x,y$ bulat, maka nilai $x+y$ adalah \dots

    \item (OSK 2015) Diketahui bilangan real positif $a$ dan $b$ memenuhi persamaan
    \begin{align*}
        a^4+a^2b^2+b^4=6 \text{ dan } a^2+ab+b^2=4
    \end{align*}
    Nilai dari $a + b$ adalah \ldots

    \item (OSK 2022) Diketahui $a,b,c,d$ bilangan real positif yang memenuhi $a>c$, $d>b$, dan 
    $$3a^2+3b^2=3c^2+3d^2=4ac+4bd.$$
    Nilai $\dfrac{12(ab+cd)}{ad+bc}=\dots$ 
\end{enumerate}

\section{Latihan Soal Eksponen}
\begin{enumerate}
    \item Carilah jumlah semua bilangan bulat positif $a$ yang memenuhi $a^{(a-1)^{(a-2)}}=a^{a^2-3a+2}$.
    
    \item Carilah jumlah seluruh solusi real $x$ yang memenuhi $(x^2+5x+5)^{x^2-10x+21}=1.$
    
    \item Jika $5^x=6^y=30^7$, berapakah nilai $\dfrac{xy}{x+y}$?
\end{enumerate}

\section{Latihan Soal Logaritma}
\begin{enumerate}  
    \item (OSK 2012) Jumlah semua bilangan bulat $x$ sehingga $^2 \log (x^2-4x-1)$ merupakan bilangan bulat adalah \dots
    
    \item (OSK 2014) Misalkan $x,y,z>1$ dan $w>0$. Jika $\log_x w = 4$, $\log_y w = 5$, dan $\log_{xyz} w = 2$, maka nilai $\log_z w$ adalah \dots 

    \item (OSK 2016) Misalkan $x,y,z$ adalah bilangan real positif yang memenuhi $$3 \log_x (3y) = 3 \log_{3x} (27z) = \log_{3x^4} (81yz) \neq 0.$$ Nilai dari $x^5y^4z$ adalah \dots
\end{enumerate}

\section{Latihan Soal Barisan dan Deret}
\begin{enumerate}
\item (OSK 2006) Diketahui $a+(a+1)+(a+2)+\dots+50=1139$. Jika $a$ bilangan real positif, maka $a=\dots$

\item (AIME 1984) Barisan $a_1,a_2,\dots,a_{98}$ memenuhi $a_{n+1}=a_n+1$ untuk $n=1,2,\dots,97$ dan mempunyai jumlah sama dengan $137$. Tentukan nilai dari $a_2+a_4+a_6+\dots+a_{98}$.

\item (AIME 2003) Diketahui $0<a<b<c<d$ adalah bilangan bulat dimana $a,b,c$ membentuk barisan aritmatika sedangkan $b,c,d$ membentuk barisan geometri. Jika $d-a=30$ maka tentukan nilai dari $a+b+c+d$.

\item (OSK 2009) Bilangan bulat positif terkecil $n$ dengan $n> 2009$ sehingga $$\sqrt{\dfrac{1^3+2^3+3^3+\dots+n^3}{n}}$$
merupakan bilangan bulat adalah \dots

\item (OSK 2015) Diketahui barisan bilangan real $a_1$, $a_2$, $a_3$, $\dots$, $a_n$, $\dots$ merupakan barisan geometri. Jika $a_1+a_4 = 20$, maka nilai minimal dari
\[a_1 + a_2 + a_3 + a_4 + a_5 + a_6\]
adalah \ldots

\item (OSK 2022) Jika $\displaystyle\sum_{k=1}^{\infty} \frac{2k + B}{3^{k+1}} = 10$, maka $B = \ldots $.

\item (OSK 2016) Diberikan barisan $\{a_n\}$ dan $\{b_n\}$ dengan $a_n = \dfrac{1}{n\sqrt{n}}$ dan $b_n = \dfrac{1}{\left(a+\frac{1}{n}\right)+\sqrt{1+\frac{1}{n}}}$, untuk setiap bilangan asli $n$. Misalkan $S_n = a_1b_1 + a_2b_2 + \ldots + a_nb_n$. Banyaknya
bilangan asli $n$ dengan $n \le 2016$ sehingga $S_n$ merupakan bilangan rasional adalah \ldots 

\item Tentukan nilai paling sederhana dari $\sqrt{2\sqrt{2\sqrt{2\sqrt{\dots}}}}$.

\item Tentukan nilai dari $\dfrac{1}{1\cdot2}+\dfrac{1}{2\cdot 3}+\dfrac{1}{3 \cdot 4}+\dots+\dfrac{1}{2020 \cdot 2021}.$

\item  Nilai paling sederhana dari $\left(1-\dfrac{1}{2^2}\right)\cdot\left(1-\dfrac{1}{3^2}\right)\cdot\left(1-\dfrac{1}{4^2}\right)\cdot\dots\cdot\left(1-\dfrac{1}{2021^2}\right)$ adalah \dots

\item Tentukan hasil dari jumlah $\dfrac{1}{\sqrt{1}+\sqrt{2}}+\dfrac{1}{\sqrt{2}+\sqrt{3}}+\dots+\dfrac{1}{\sqrt{99}+\sqrt{100}}.$

\item Nilai $x$ yang memenuhi persamaan
$$\sqrt{x\sqrt{x\sqrt{x\sqrt{\dots}}}}=\sqrt{4x+\sqrt{4x+\sqrt{4x+\sqrt{\dots}}}}$$
adalah \dots

\item Hitunglah nilai paling sederhana dari
$$6-\dfrac{5}{3+\dfrac{4}{3+\dfrac{4}{3+\dfrac{4}{3+\dfrac{4}{\dots}}}}}$$
\end{enumerate}

\section{Latihan Soal Fungsi}
\begin{enumerate}
        \item (OSK 2015) Jika $(f \circ g)(x) = \frac{7x + 3}{5x - 9}$ dan $g(x) = 2x - 4$, maka nilai $f(2)$ adalah \ldots

        \item Jika $f:\RR-\{0\} \rightarrow \RR$ adalah fungsi yang memenuhi $f(x)+2f\left(\frac{1}{x}\right)=3x$ untuk setiap bilangan real $x \neq 0$, maka tentukan nilai $f(2024).$

        \item (OSP 2004) Misalkan $f$ sebuah fungsi yang memenuhi $f(x)f(y)-f(xy)=x+y,$ untuk setiap bilangan bulat $x$ dan $y$. Berapakah nilai $f(2004)$?
        
        \item (OSK 2011) Misalkan $f$ suatu fungsi yang memenuhi $f(xy) = \dfrac{f(x)}{y}$ untuk semua bilangan real positif $x$ dan $y$. Jika $f(100)=3$ maka $f(10)$ adalah \dots
        
        \item (OSP 2009) Suatu fungsi $f:\ZZ \rightarrow \QQ$ mempunyai sifat $f(x+1)=\dfrac{1+f(x)}{1-f(x)}$ untuk setiap $x \in \ZZ$. Jika $f(2)=2$, maka nilai fungsi $f(2009)$ adalah \dots

        \item (OSK 2022) Misalkan $f(x) = a^{2x} + 300$. Jika
        $$f(20) + f^{-1}(22) = f^{-1}(20) + f(22)$$,
        maka $f(1) = \dots$.

\end{enumerate}

\section{Latihan Soal Polinomial}
\begin{enumerate}

\item Jika $P(x)$ dibagi $x^2-x$ dan $x^2+x$ berturut-turut akan bersisa $5x+1$ dan $3x+1$, maka bila $P(x)$ dibagi $x^2-1$ sisanya adalah \dots

\item (OSP 2006) Jika $(x-1)^2$ membagi $ax^4+bx^3+1$, maka $ab=\dots$

\item Diketahui suatu polinomial $P(x)$ memenuhi $P(k)=\dfrac{k}{k+1}$ untuk $k=1,2,3,\dots,2020$. Jika $P(0)=1$, nilai $P(2022)=\dots$

\item (OSK 2010) Polinom $P(x)=x^3-x^2+x-2$ mempunyai tiga pembuat nol yaitu $a,b,$ dan $c$. Nilai dari $a^3+b^3+c^3$ adalah \dots

\item (OSP 2010) Persamaan kuadrat $x^2-px-2p=0$ mempunyai dua akar real $a$ dan $b$. Jika $a^3+b^3=16$, maka hasil jumlah semua nilai $p$ yang memenuhi adalah \dots 

\item (OSP 2010) Diberikan polinomial $P(x)=x^4+ax^3+bx^2+cx+d$ dengan $a,b,c,$ dan $d$ konstanta. Jika $P(1)=10$, $P(2)=20$, dan  $P(3)=30$, maka nilai
$$\dfrac{P(12)+P(-8)}{10}=\dots$$

\item (OSK 2015) Diketahui $a$, $b$, $c$ akar-akar dari persamaan $x^3 - 5x^2 - 9x + 10 = 0$. Jika suku banyak $P(x) = Ax^3 + Bx^2 + Cx - 2015$ memenuhi $P(a) = b + c$, $P(b) = a + c$ dan $P(c) = a + b$ maka nilai dari $A + B + C$ adalah \ldots

\end{enumerate}

\section{Latihan Soal Ketaksamaan}
\begin{enumerate}
    \item (Nesbitt's Inequality) Untuk bilangan real positif $a,b,c$ tentukan nilai minimum dari $$\dfrac{a}{b+c}+\dfrac{b}{c+a}+\dfrac{c}{a+b}.$$
    
    \item Tentukan nilai minimum dari $8x^4+y^2$ untuk bilangan real positif $x$ dan $y$ yang memenuhi $x^4y=\dfrac{1}{\sqrt{2}}$.
    
    \item Nilai minimum dari $$\dfrac{1}{w}+\dfrac{1}{x}+\dfrac{1}{y}+\dfrac{1}{z}$$
    untuk bilangan real positif $w,x,y,z$ yang memenuhi $w+x+y+z=32$ adalah \dots
    
    \item Jika $x^2+y^2+z^2=1$, nilai maksimum dari $x+2y+3z$ adalah \dots
    
    \item Diberikan $a+b+c=1$ dan $a,b,c>0$, carilah nilai minimum dari $a^2+2b^2+c^2$.
    
    \item (OSK 2014) Untuk $0 < x < \pi$, nilai minimum dari $\dfrac{16 \sin^2 x + 9}{\sin x}$ adalah \dots
    
    \item (OSK 2017) Misalkan $a,b,c$ bilangan real positif yang memenuhi $a+b+c=1$. Nilai minimum dari $\dfrac{a+b}{abc}$ adalah \dots
    
    \item (OSK 2017) Pada segitiga $ABC$ titik $K$ dan $L$ berturut-turut adalah titik tengah $AB$ dan $AC$. Jika $CK$ dan $BL$ saling tegak lurus, maka nilai minimum dari $\cot B + \cot C$ adalah \dots
\end{enumerate}

\section{Latihan Soal Floor dan Ceiling}
\begin{enumerate}
    \item (Modifikasi JBMO 2021) Carilah seluruh penyelesaian dari persamaan $2\cdot \lfloor{\frac{1}{2x}}\rfloor - 7 = 9(1 - 8x)$.

    \item (OSK 2013) Misalkan $\floor{x}$ menyatakan bilangan bulat terbesar yang lebih kecil atau sama dengan $x$ dan $\ceiling{x}$ menyatakan bilangan bulat terkecil yang lebih besar atau sama dengan $x$. Tentukan semua $x$ yang memenuhi $\floor{x}$ + $\ceiling{x}$ = 5.
    
    \item (OSK 2016) Banyaknya bilangan asli $n \in \{1,2,3,\dots,1000\}$ sehingga terdapat bilangan real positif $x$ yang memenuhi $x^2+\floor{x}^2=n$ adalah \dots
    
    \item (OSK 2018) Untuk setiap bilangan real $z$, $\lfloor z \rfloor$ menyatakan bilangan bulat terbesar yang lebih kecil dari atau sama dengan $z$. Jika diketahui $\lfloor x \rfloor + \lfloor y \rfloor + y = 43.8$ dan $x + y - \lfloor x \rfloor = 18.4$. Nilai $10(x + y)$ adalah...
    
    \item Let $[x]$ denote the largest integer not exceeding $x$. For example, $[2.1]=2$, $[4]=4$ and $[5.7]=5$. How many positive integers $n$ satisfy the equation $\left[\frac{n}{5}\right]=\frac{n}{6}$.

    \item (OSK 2019) Untuk sebarang bilangan real $x$, simbol $\lfloor x \rfloor$ menyatakan bilangan bulat terbesar yang tidak lebih besar daripada $x$, sedangkan $\lceil x \rceil$ menyatakan bilangan bulat terkecil yang tidak lebih kecil dibanding $x$. Interval $[a, b)$ adalah himpunan semua bilangan real $x$ yang memenuhi
    $$\lfloor 2x \rfloor^2 = \lceil x \rceil + 7$$.
    Nilai $a \cdot b$ adalah ...

    \item (OSK 2021) Jika $a > 1$ suatu bilangan asli sehingga hasil penjumlahan semua bilangan riil $x$ yang memenuhi persamaan
    $$\lfloor x \rfloor^2 - 2ax + a = 0$$
    adalah $p$, maka $a$ adalah ...

    \item Carilah banyaknya bilangan real $x$ yang memenuhi
    $$\floor{x^2}=4x+3.$$

    \item Jika $x$ adalah suatu bilangan real yang memenuhi persamaan berikut
    $$\floor{x}+\floor{2x}+\floor{3x}+\floor{4x}=2024.$$

    Tentukan nilai dari $\floor{6x}$.

    \item Jumlah dari semua solusi bilangan real $x$ dari persamaan
    $$2\floor{x}^2+3\{x\}^2=\frac{7}{4}x\floor{x}$$
    dapat dinyatakan dalam bentuk $\frac{p}{q}$ dimana $p$ dan $q$ merupakan bilangan bulat yang saling relatif prima. Tentukan nilai dari $10p+q$.
\end{enumerate}

\end{document}